\section{Roteiro Prático 05 — Assincronismo, Fetch API \& Backend RESTful em Flask (SQLite)}

% =====================================================
% OBJETIVOS DA ATIVIDADE
% =====================================================

\begin{boxObjetivos}
\textbf{Objetivos de Aprendizagem (Módulo 3 — Assincronismo \& Integração de Serviços Web)}

Ao final desta prática orientada em laboratório, o estudante será capaz de:

\begin{itemize}[leftmargin=*]
\item Compreender a transição do armazenamento isolado no navegador (\texttt{localStorage}) para a \textbf{arquitetura Cliente-Servidor real}, onde os dados são centralizados e compartilhados em rede;
\item Configurar e executar um microbackend RESTful em \textbf{Python Flask}, compreendendo o papel dos protocolos HTTP, rotas, verbos (\texttt{GET}, \texttt{POST}, \texttt{PUT}, \texttt{DELETE}) e códigos de status (\texttt{200}, \texttt{201}, \texttt{400}, \texttt{404}, \texttt{500});
\item Integrar um banco de dados relacional embutido (\textbf{SQLite 3}) via arquivo persistente (\texttt{vitrine.db}), eliminando a necessidade de instalar servidores externos de banco;
\item Compreender o mecanismo de segurança \textbf{CORS} (\textit{Cross-Origin Resource Sharing}) e por que ele bloqueia requisições entre portas distintas (ex: frontend na porta \texttt{5500} e API na porta \texttt{5000});
\item Refatorar a camada de serviços (\texttt{vitrineService.js}) para operar com a moderna \textbf{Fetch API} do JavaScript nativo utilizando a sintaxe \textbf{\texttt{async/await}} e tratamento robusto de exceções com blocos \textbf{\texttt{try/catch}};
\item Implementar o padrão profissional de \textbf{Contingência Resiliente (Fallback)}: se a API do Flask estiver ativa (no laboratório), o sistema opera conectado ao banco relacional; se a página for executada no GitHub Pages ou offline, a aplicação chaveia automaticamente para o \texttt{localStorage} sem quebrar;
\item Implementar estados visuais reativos no frontend com \textbf{Spinners de Carregamento} e \textbf{Badges dinâmicos de conectividade} com a API.
\end{itemize}
\end{boxObjetivos}

% =====================================================
% INTRODUÇÃO E EVOLUÇÃO ARQUITETURAL
% =====================================================

\subsection{Introdução: Da Vitrine Isolada à Aplicação em Rede}

No \textbf{Roteiro Prático 04}, demos um grande salto organizando o frontend na arquitetura modular com ES6 Modules (\texttt{services/}, \texttt{views/}, \texttt{utils/}). No entanto, todo o armazenamento ocorria exclusivamente dentro do navegador de cada máquina através da \texttt{Web Storage API} (\texttt{localStorage}).

Na prática profissional, se um empreendedor em Garopaba cadastrar um novo serviço em seu celular, esse serviço precisa ficar imediatamente disponível para qualquer cliente que abrir a página em outro computador. Isso exige que os dados saiam do navegador e sejam gravados em um **servidor central com banco de dados**.

Neste \textbf{Roteiro Prático 05}, evoluiremos a mesma base de código da \textbf{Vitrine Comunitária}, conectando o frontend modular a uma API RESTful leve construída em **Flask** com banco de dados relacional **SQLite**.

% =====================================================
% DIAGRAMA DA ARQUITETURA CLIENTE-SERVIDOR
% =====================================================

\subsection{Visão Geral da Arquitetura Cliente-Servidor}

\begin{center}
\begin{tikzpicture}[
    node distance=1.5cm and 2.2cm,
    caixa/.style={rectangle, draw=azulTitulo, rounded corners=5pt, fill=azulTitulo!6, line width=1pt, align=center, inner sep=8pt, font=\small},
    nuvem/.style={rectangle, draw=ifscGreen, rounded corners=8pt, fill=ifscGreen!10, line width=1.2pt, align=center, inner sep=10pt, font=\small\bfseries},
    seta/.style={-{Latex[length=3mm]}, line width=1.2pt}
]

% Nós
\node[caixa] (cliente) {\textbf{Cliente Web (Navegador)}\\Bootstrap 5 + ES6 Modules\\\texttt{Live Server (Porta 5500)}};

\node[nuvem, right=of cliente] (api) {\textbf{Backend RESTful (Flask)}\\Microframework Python\\\texttt{http://127.0.0.1:5000}};

\node[caixa, right=of api] (banco) {\textbf{Banco Relacional}\\\textbf{SQLite 3 em Arquivo}\\\texttt{vitrine.db}};

% Conexões
\draw[seta, color=azulTitulo] (cliente.north east) to[bend left=20] node[midway, above, font=\footnotesize] {Requisição HTTP: \texttt{fetch('/api/servicos')}} (api.north west);

\draw[seta, color=ifscGreen] (api.south west) to[bend left=20] node[midway, below, font=\footnotesize] {Resposta JSON: \texttt{[\{id: 1, nome: "..."\}]}} (cliente.south east);

\draw[seta, color=pyKeyword] (api.east) -- node[midway, above, font=\footnotesize] {SQL Query} (banco.west);
\draw[seta, color=pyKeyword] (banco.south west) to[bend left=15] node[midway, below, font=\footnotesize] {Cursor / Rows} (api.south east);

\end{tikzpicture}
\end{center}

% =====================================================
% ESTRUTURA DE DIRETÓRIOS DO PROJETO
% =====================================================

\subsection{Estrutura de Pastas do Projeto Integrado}

Para manter o desacoplamento saudável entre a camada de backend e a camada de frontend, organizaremos o projeto na pasta \texttt{roteiro-5/exercicio-roteiro5/} com a seguinte estrutura:

\begin{boxCodigo}
\textbf{Árvore de Diretórios do Projeto Integrado (Full-Stack Didático):}

\fontsize{8.5pt}{11pt}\selectfont
\begin{verbatim}
roteiro-5/exercicio-roteiro5/
|-- backend/
|   |-- requirements.txt   <-- Dependências Python (Flask, flask-cors)
|   |-- app.py             <-- API RESTful com rotas e conexão SQLite
|   `-- vitrine.db         <-- Banco SQLite (gerado automaticamente no 1º run)
`-- frontend/
    |-- index.html         <-- Interface com status de conexão e spinner
    |-- styles.css         <-- Estilos complementares
    `-- js/
        |-- main.js        <-- Orquestrador assíncrono de eventos
        |-- services/
        |   `-- vitrineService.js <-- Consumo Fetch API com fallback
        |-- views/
        |   `-- vitrineView.js    <-- Renderizador de cards, spinner e toasts
        `-- utils/
            `-- formatters.js     <-- Formatações (100% inalterado!)
\end{verbatim}
\end{boxCodigo}

\vspace{0.2cm}
\begin{tcolorbox}[colback=ifscGreen!10, colframe=ifscGreen, title=\textbf{Princípio do Desacoplamento Arquitetural}, arc=3pt]
Observe atentamente que a camada utilitária (\texttt{formatters.js}) e grande parte da camada de visão (\texttt{vitrineView.js}) **permanecem exatamente as mesmas do Roteiro 04**. Toda a migração de tecnologia ocorre de forma isolada no \texttt{vitrineService.js}. Este é o maior benefício do padrão arquitetural que aprendemos!
\end{tcolorbox}

% =====================================================
% FUNDAMENTAÇÃO TEÓRICA
% =====================================================

\subsection{Fundamentação Teórica: Conceitos-Chave}

\subsubsection{1. O Ciclo de Vida HTTP e Verbos REST}
Uma API RESTful utiliza o protocolo HTTP padrão para realizar operações sobre recursos:
\begin{itemize}[leftmargin=*]
  \item \textbf{\texttt{GET /api/servicos}:} Solicita a leitura da lista de serviços cadastrados (\textit{Read}). Não altera o estado do servidor. Retorna código de status \textbf{\texttt{200 OK}}.
  \item \textbf{\texttt{POST /api/servicos}:} Envia um objeto JSON no corpo (\textit{Body}) da requisição para cadastrar um novo serviço (\textit{Create}). Retorna código \textbf{\texttt{201 Created}}.
  \item \textbf{\texttt{PUT /api/servicos/<id>}:} Atualiza os dados de um registro existente pelo seu identificador (\textit{Update}). Retorna código \textbf{\texttt{200 OK}}.
  \item \textbf{\texttt{DELETE /api/servicos/<id>}:} Remove permanentemente um serviço do banco (\textit{Delete}). Retorna código \textbf{\texttt{200 OK}}.
\end{itemize}

\subsubsection{2. A Política de Mesma Origem e o CORS}
Por questões de segurança, os navegadores web impedem por padrão que um script executado em uma origem (ex: \texttt{http://127.0.0.1:5500} no Live Server) leia dados de outra origem (ex: \texttt{http://127.0.0.1:5000} no Flask). Essa política de bloqueio é chamada de **SOP** (\textit{Same-Origin Policy}).

Para autorizar que o nosso frontend acesse a API, o servidor Flask precisa enviar nos cabeçalhos de resposta a diretiva \textbf{CORS} (\textit{Cross-Origin Resource Sharing}):
\begin{center}
\texttt{Access-Control-Allow-Origin: *}
\end{center}
No Python, essa autorização é resolvida de maneira elegante através da extensão \textbf{\texttt{flask-cors}}.

\subsubsection{3. Assincronismo com Promises e async/await}
No JavaScript, operações de rede (como buscar dados em outro computador) levam alguns milissegundos ou segundos para responder. Se o código executasse de forma síncrona (\textit{bloqueante}), o navegador congelaria completamente enquanto esperasse a resposta.

Por isso, o JavaScript utiliza **Promises** (promessas de valores futuros) manipuladas pelas palavras-chave \textbf{\texttt{async}} e \textbf{\texttt{await}}:
\begin{boxCodigoJS}
// A palavra 'async' declara que a função lida com operações assíncronas
export async function obterDados() {
  try {
    // A palavra 'await' pausa a execução DESTA função até a resposta chegar,
    // permitindo que o navegador continue fluido e responsivo para o usuário.
    const resposta = await fetch('http://127.0.0.1:5000/api/servicos');
    
    // Converte o fluxo de texto HTTP em um objeto JavaScript
    const dados = await resposta.json();
    return dados;
  } catch (erro) {
    console.error('Falha de rede:', erro);
  }
}
\end{boxCodigoJS}

% =====================================================
% FASES DE IMPLEMENTAÇÃO GUIADAS (ZERO-ELLIPSIS RULE)
% =====================================================

\newpage
\subsection{Implementação Passo a Passo Guiada}

Seguiremos a sequência didática em 6 fases de implementação estruturadas. Todo o código fornecido é **100\% completo, sem reticências ou omissões**.

% --- FASE 1 ---
\begin{boxAtividade}[{Fase 1/6: Configuração do Backend em Flask \& SQLite (backend/)}]
\textbf{Responsabilidade Arquitetural:} Criar o ambiente isolado Python (virtualenv), instalar o microframework Flask com suporte a CORS e implementar a API RESTful com conexão ao SQLite em arquivo (\texttt{vitrine.db}).
\end{boxAtividade}

\textbf{Passo 1.1: Criar o arquivo de dependências (\texttt{backend/requirements.txt})}
\begin{boxCodigoTerminal}
\fontsize{8.5pt}{11pt}\selectfont
\begin{verbatim}
Flask>=3.0.0
flask-cors>=4.0.0
\end{verbatim}
\end{boxCodigoTerminal}

\textbf{Passo 1.2: No terminal do VS Code, criar o ambiente virtual e instalar dependências:}
\begin{boxCodigoTerminal}
\fontsize{8.5pt}{11pt}\selectfont
\begin{verbatim}
# Acesse a pasta do backend
cd roteiro-5/exercicio-roteiro5/backend

# Cria o ambiente virtual isolado
python3 -m venv venv

# Ativa o ambiente virtual (Linux/Mac)
source venv/bin/activate
# No Windows PowerShell utilize: .\venv\Scripts\Activate.ps1

# Instala o Flask e o CORS
pip install -r requirements.txt
\end{verbatim}
\end{boxCodigoTerminal}

\textbf{Passo 1.3: Criar o arquivo completo do servidor (\texttt{backend/app.py}):}
\begin{boxCodigoPY}
import os
import sqlite3
from datetime import datetime
from flask import Flask, jsonify, request
from flask_cors import CORS

app = Flask(__name__)
# Habilita CORS para permitir que o Live Server (porta 5500) acesse o Flask (5000)
CORS(app)

DB_PATH = os.path.join(os.path.dirname(__file__), 'vitrine.db')

def obter_conexao():
    """Conecta ao SQLite retornando linhas como dicionários."""
    conn = sqlite3.connect(DB_PATH)
    conn.row_factory = sqlite3.Row
    return conn

def inicializar_banco():
    """Cria a tabela servicos e popula com os 3 registros modelo de Garopaba se vazia."""
    conn = obter_conexao()
    cursor = conn.cursor()
    cursor.execute('''
        CREATE TABLE IF NOT EXISTS servicos (
            id INTEGER PRIMARY KEY AUTOINCREMENT,
            nome TEXT NOT NULL,
            categoria TEXT NOT NULL,
            bairro TEXT NOT NULL,
            precoBase REAL NOT NULL,
            telefone TEXT NOT NULL,
            descricao TEXT,
            dataCadastro TEXT
        )
    ''')

    cursor.execute('SELECT COUNT(*) AS total FROM servicos')
    if cursor.fetchone()['total'] == 0:
        dados_semente = [
            (
                'Maré Alta Artesanatos & Cerâmicas', 'Artesanato', 'Centro Histórico',
                35.00, '48991234567',
                'Peças artesanais e utilitárias modeladas à mão com argila local e conchas de Garopaba.',
                datetime.now().isoformat()
            ),
            (
                'Garopaba Web & Design Studio', 'Tecnologia', 'Ferrugem',
                150.00, '48998765432',
                'Criação de websites profissionais responsivos, cardápios digitais e suporte para comércio local.',
                datetime.now().isoformat()
            ),
            (
                'Pescado Fresco do Zequinha', 'Alimentação', 'Canto das Canoas',
                42.00, '48984561234',
                'Peixes frescos e frutos do mar da pesca artesanal diária entregues com higiene e pontualidade.',
                datetime.now().isoformat()
            )
        ]
        cursor.executemany('''
            INSERT INTO servicos (nome, categoria, bairro, precoBase, telefone, descricao, dataCadastro)
            VALUES (?, ?, ?, ?, ?, ?, ?)
        ''', dados_semente)
        conn.commit()
    conn.close()

@app.route('/api/status', methods=['GET'])
def verificar_status():
    """Diagnóstico de conectividade com a API."""
    return jsonify({
        'status': 'online',
        'mensagem': 'API REST da Vitrine de Garopaba operacional.',
        'banco': 'SQLite 3'
    }), 200

@app.route('/api/servicos', methods=['GET'])
def listar_servicos():
    """READ: Retorna a lista de todos os serviços ordenados pelo ID decrescente."""
    conn = obter_conexao()
    cursor = conn.cursor()
    cursor.execute('SELECT * FROM servicos ORDER BY id DESC')
    linhas = cursor.fetchall()
    conn.close()
    return jsonify([dict(l) for l in linhas]), 200

@app.route('/api/servicos', methods=['POST'])
def cadastrar_servico():
    """CREATE: Grava um novo serviço enviado em formato JSON no SQLite."""
    dados = request.get_json()
    if not dados:
        return jsonify({'erro': 'Payload inválido.'}), 400

    conn = obter_conexao()
    cursor = conn.cursor()
    data_atual = datetime.now().isoformat()

    cursor.execute('''
        INSERT INTO servicos (nome, categoria, bairro, precoBase, telefone, descricao, dataCadastro)
        VALUES (?, ?, ?, ?, ?, ?, ?)
    ''', (
        dados['nome'].strip(),
        dados['categoria'].strip(),
        dados['bairro'].strip(),
        float(dados['precoBase']),
        dados['telefone'].strip(),
        dados.get('descricao', '').strip(),
        data_atual
    ))
    conn.commit()
    novo_id = cursor.lastrowid
    conn.close()

    return jsonify({
        'id': novo_id,
        'nome': dados['nome'].strip(),
        'categoria': dados['categoria'].strip(),
        'bairro': dados['bairro'].strip(),
        'precoBase': float(dados['precoBase']),
        'telefone': dados['telefone'].strip(),
        'descricao': dados.get('descricao', '').strip(),
        'dataCadastro': data_atual
    }), 201

@app.route('/api/servicos/<int:id>', methods=['DELETE'])
def excluir_servico(id):
    """DELETE: Exclui um serviço pelo ID."""
    conn = obter_conexao()
    cursor = conn.cursor()
    cursor.execute('DELETE FROM servicos WHERE id = ?', (id,))
    conn.commit()
    conn.close()
    return jsonify({'mensagem': f'Serviço {id} removido.'}), 200

if __name__ == '__main__':
    inicializar_banco()
    print('Servidor Flask ativo na porta 5000: http://127.0.0.1:5000')
    app.run(host='127.0.0.1', port=5000, debug=True)
\end{boxCodigoPY}

% --- FASE 2 ---
\newpage
\begin{boxAtividade}[{Fase 2/6: Executando e Testando a API RESTful}]
\textbf{Responsabilidade Arquitetural:} Iniciar o servidor Flask na porta 5000 e validar os endpoints de saúde (\texttt{/api/status}) e listagem (\texttt{/api/servicos}) diretamente no terminal e navegador antes da integração.
\end{boxAtividade}

\textbf{Passo 2.1: Iniciar o servidor no terminal:}
\begin{boxCodigoTerminal}
\fontsize{8.5pt}{11pt}\selectfont
\begin{verbatim}
python app.py
\end{verbatim}
\end{boxCodigoTerminal}

\textbf{Passo 2.2: Testar no navegador ou via cURL em outra aba do terminal:}
\begin{boxCodigoTerminal}
\fontsize{8.5pt}{11pt}\selectfont
\begin{verbatim}
curl http://127.0.0.1:5000/api/status
\end{verbatim}
\end{boxCodigoTerminal}

\textbf{Resposta esperada do servidor (JSON com status 200):}
\begin{boxCodigoTerminal}
\fontsize{8.5pt}{11pt}\selectfont
\begin{verbatim}
{"banco":"SQLite 3","mensagem":"API REST da Vitrine de Garopaba operacional.","status":"online"}
\end{verbatim}
\end{boxCodigoTerminal}

\begin{boxTeste}{Checkpoint de Validação da Fase 2}
Abra no navegador \url{http://127.0.0.1:5000/api/servicos}. Você deverá ver a lista JSON com os 3 registros semente cadastrados automaticamente no arquivo \texttt{vitrine.db}.
\end{boxTeste}

% --- FASE 3 ---
\begin{boxAtividade}[{Fase 3/6: A Revolução Assíncrona no vitrineService.js}]
\textbf{Responsabilidade Arquitetural:} Refatorar os métodos de leitura e escrita para utilizarem requisições HTTP assíncronas com \texttt{fetch()} e \texttt{async/await}, mantendo uma camada de contingência transparente com \texttt{localStorage}.
\end{boxAtividade}

\begin{boxCodigoJS}
// js/services/vitrineService.js
const API_BASE_URL = 'http://127.0.0.1:5000/api';
const STORAGE_KEY = 'garopaba_vitrine_servicos';

let apiConectada = false;

export const DADOS_INICIAIS = [
  {
    id: '1',
    nome: 'Maré Alta Artesanatos & Cerâmicas',
    categoria: 'Artesanato',
    bairro: 'Centro Histórico',
    precoBase: 35.00,
    telefone: '48991234567',
    descricao: 'Peças artesanais e utilitárias modeladas à mão com argila local e conchas de Garopaba.'
  },
  {
    id: '2',
    nome: 'Garopaba Web & Design Studio',
    categoria: 'Tecnologia',
    bairro: 'Ferrugem',
    precoBase: 150.00,
    telefone: '48998765432',
    descricao: 'Criação de websites profissionais responsivos, cardápios digitais e suporte para comércio local.'
  },
  {
    id: '3',
    nome: 'Pescado Fresco do Zequinha',
    categoria: 'Alimentação',
    bairro: 'Canto das Canoas',
    precoBase: 42.00,
    telefone: '48984561234',
    descricao: 'Peixes frescos e frutos do mar da pesca artesanal diária entregues com higiene e pontualidade.'
  }
];

export function estaConectadoNaApi() {
  return apiConectada;
}

export async function testarConexaoApi() {
  try {
    const resposta = await fetch(`${API_BASE_URL}/status`, {
      method: 'GET',
      headers: { 'Accept': 'application/json' },
      signal: AbortSignal.timeout(2000)
    });
    apiConectada = resposta.ok;
    return apiConectada;
  } catch (erro) {
    apiConectada = false;
    return false;
  }
}

// Rotinas Locais de Contingência (Fallback)
function obterServicosLocal() {
  const dados = localStorage.getItem(STORAGE_KEY);
  if (!dados) {
    localStorage.setItem(STORAGE_KEY, JSON.stringify(DADOS_INICIAIS));
    return DADOS_INICIAIS;
  }
  try { return JSON.parse(dados); } catch (e) { return []; }
}

function salvarServicoLocal(novoServico) {
  const servicos = obterServicosLocal();
  const servicoCompleto = {
    id: Date.now().toString(),
    dataCadastro: new Date().toISOString(),
    ...novoServico
  };
  servicos.unshift(servicoCompleto);
  localStorage.setItem(STORAGE_KEY, JSON.stringify(servicos));
  return servicoCompleto;
}

function removerServicoLocal(id) {
  const servicos = obterServicosLocal();
  const filtrados = servicos.filter(s => s.id.toString() !== id.toString());
  localStorage.setItem(STORAGE_KEY, JSON.stringify(filtrados));
  return filtrados;
}

// Operações Assíncronas com Fetch API
export async function obterServicos() {
  try {
    const resposta = await fetch(`${API_BASE_URL}/servicos`, {
      method: 'GET',
      headers: { 'Accept': 'application/json' },
      signal: AbortSignal.timeout(3000)
    });
    if (!resposta.ok) throw new Error(`HTTP ${resposta.status}`);
    const dados = await resposta.json();
    apiConectada = true;
    localStorage.setItem(STORAGE_KEY, JSON.stringify(dados));
    return dados;
  } catch (erro) {
    apiConectada = false;
    console.warn('API Flask offline. Ativando contingência LocalStorage:', erro.message);
    return obterServicosLocal();
  }
}

export async function salvarServico(novoServico) {
  try {
    const resposta = await fetch(`${API_BASE_URL}/servicos`, {
      method: 'POST',
      headers: {
        'Content-Type': 'application/json',
        'Accept': 'application/json'
      },
      body: JSON.stringify(novoServico),
      signal: AbortSignal.timeout(4000)
    });
    if (!resposta.ok) throw new Error(`HTTP ${resposta.status}`);
    const criado = await resposta.json();
    apiConectada = true;
    const servicos = obterServicosLocal();
    servicos.unshift(criado);
    localStorage.setItem(STORAGE_KEY, JSON.stringify(servicos));
    return criado;
  } catch (erro) {
    apiConectada = false;
    console.warn('Falha no envio à API. Salvando localmente:', erro.message);
    return salvarServicoLocal(novoServico);
  }
}

export async function removerServico(id) {
  try {
    const resposta = await fetch(`${API_BASE_URL}/servicos/${id}`, {
      method: 'DELETE',
      headers: { 'Accept': 'application/json' },
      signal: AbortSignal.timeout(4000)
    });
    if (!resposta.ok) throw new Error(`HTTP ${resposta.status}`);
    apiConectada = true;
    removerServicoLocal(id);
    return true;
  } catch (erro) {
    apiConectada = false;
    console.warn('Falha na exclusão remota. Removendo localmente:', erro.message);
    removerServicoLocal(id);
    return true;
  }
}
\end{boxCodigoJS}

% --- FASE 4 ---
\newpage
\begin{boxAtividade}[{Fase 4/6: Camada de Visão \& Feedbacks de Rede (vitrineView.js)}]
\textbf{Responsabilidade Arquitetural:} Adicionar componentes visuais de carregamento (\texttt{exibirSpinner}) e indicador de conectividade em tempo real com a API no topo da página.
\end{boxAtividade}

\begin{boxCodigoJS}
// js/views/vitrineView.js
import { formatarMoeda, formatarTelefone, escapeHtml } from '../utils/formatters.js';

const CORES_CATEGORIA = {
  'Alimentação': 'success',
  'Tecnologia': 'info',
  'Artesanato': 'warning',
  'Serviços Gerais': 'primary',
  'Turismo': 'secondary'
};

export function criarCardHtml(s) {
  const corBadge = CORES_CATEGORIA[s.categoria] || 'primary';
  const telNumeros = (s.telefone || '').replace(/\D/g, '');

  return `
    <div class="col-md-6 col-lg-4 mb-4">
      <div class="card h-100 shadow-sm border-0 vitrine-card rounded-4 overflow-hidden">
        <div class="card-body p-4 d-flex flex-column">
          <div class="d-flex justify-content-between align-items-center mb-2">
            <span class="badge bg-${corBadge} px-3 py-2 rounded-pill font-outfit">
              <i class="bi bi-tag-fill me-1"></i>${escapeHtml(s.categoria)}
            </span>
            <span class="fw-bold text-success font-outfit fs-5">
              ${formatarMoeda(s.precoBase)}
            </span>
          </div>

          <h5 class="card-title fw-bold text-dark mt-2 mb-1">${escapeHtml(s.nome)}</h5>
          <h6 class="text-muted small mb-3">
            <i class="bi bi-geo-alt-fill text-danger me-1"></i>${escapeHtml(s.bairro)}
          </h6>
          <p class="card-text text-secondary small flex-grow-1" style="line-height: 1.5;">
            ${escapeHtml(s.descricao)}
          </p>

          <hr class="my-3 text-muted opacity-25">

          <div class="d-flex justify-content-between align-items-center mt-auto">
            <a href="https://wa.me/55${telNumeros}" target="_blank" class="btn btn-sm btn-outline-success rounded-pill px-3 fw-bold">
              <i class="bi bi-whatsapp me-1"></i>${formatarTelefone(s.telefone)}
            </a>
            <div class="d-flex gap-1">
              <button class="btn btn-sm btn-outline-danger btn-excluir rounded-pill px-2" data-id="${s.id}" title="Remover da vitrine">
                <i class="bi bi-trash"></i>
              </button>
            </div>
          </div>
        </div>
      </div>
    </div>
  `;
}

export function renderizarCards(servicos, containerElement) {
  if (!containerElement) return;
  if (!servicos || servicos.length === 0) {
    containerElement.innerHTML = `
      <div class="col-12 py-5 text-center text-muted">
        <i class="bi bi-inbox fs-1 d-block mb-3 text-secondary opacity-50"></i>
        <h5 class="fw-bold">Nenhum serviço encontrado</h5>
        <p class="small">Cadastre um novo serviço ou altere o filtro de categoria.</p>
      </div>
    `;
    return;
  }
  containerElement.innerHTML = servicos.map(s => criarCardHtml(s)).join('');
}

export function exibirSpinner(containerElement) {
  if (!containerElement) return;
  containerElement.innerHTML = `
    <div class="col-12 py-5 text-center text-primary">
      <div class="spinner-border text-primary" role="status" style="width: 3rem; height: 3rem;">
        <span class="visually-hidden">Carregando dados da API...</span>
      </div>
      <p class="mt-3 text-muted small fw-semibold">Consultando serviços no backend Flask (SQLite)...</p>
    </div>
  `;
}

export function atualizarStatusConexao(isOnline, statusBadgeEl) {
  if (!statusBadgeEl) return;
  if (isOnline) {
    statusBadgeEl.className = 'badge bg-success-subtle text-success border border-success-subtle px-3 py-2 rounded-pill fw-semibold';
    statusBadgeEl.innerHTML = '<i class="bi bi-cloud-check-fill me-1"></i> Backend Online (Flask + SQLite)';
  } else {
    statusBadgeEl.className = 'badge bg-warning-subtle text-warning border border-warning-subtle px-3 py-2 rounded-pill fw-semibold';
    statusBadgeEl.innerHTML = '<i class="bi bi-hdd-network-fill me-1"></i> Modo Offline (LocalStorage)';
  }
}

export function exibirToast(mensagem, tipo = 'success') {
  const toastEl = document.getElementById('toastNotificacao');
  const toastMsgEl = document.getElementById('toastMensagem');
  if (!toastEl || !toastMsgEl) return;
  toastEl.className = `toast align-items-center text-white bg-${tipo} border-0 shadow-lg`;
  toastMsgEl.innerText = mensagem;
  const bsToast = bootstrap.Toast.getOrCreateInstance(toastEl, { delay: 3500 });
  bsToast.show();
}
\end{boxCodigoJS}

% --- FASE 5 ---
\newpage
\begin{boxAtividade}[{Fase 5/6: Orquestrador Assíncrono da Aplicação (main.js)}]
\textbf{Responsabilidade Arquitetural:} Coordenar os fluxos de eventos da interface utilizando \texttt{await} para aguardar a persistência no banco remoto antes de fechar modais e atualizar a vitrine.
\end{boxAtividade}

\begin{boxCodigoJS}
// js/main.js
import {
  obterServicos,
  salvarServico,
  removerServico,
  testarConexaoApi,
  estaConectadoNaApi
} from './services/vitrineService.js';

import {
  renderizarCards,
  exibirSpinner,
  atualizarStatusConexao,
  exibirToast
} from './views/vitrineView.js';

document.addEventListener('DOMContentLoaded', async () => {
  const form = document.getElementById('formCadastro');
  const container = document.getElementById('vitrineContainer');
  const filtroCategoria = document.getElementById('filtroCategoria');
  const contadorEl = document.getElementById('totalServicosBadge');
  const statusApiBadge = document.getElementById('statusApiBadge');
  const modalEl = document.getElementById('modalCadastro');
  const btnSalvarTexto = document.getElementById('btnSalvarTexto');
  const btnSalvar = document.getElementById('btnSalvar');

  let cacheServicosAtuais = [];

  async function atualizarInterface() {
    exibirSpinner(container);
    const conectado = await testarConexaoApi();
    atualizarStatusConexao(conectado, statusApiBadge);

    cacheServicosAtuais = await obterServicos();
    const categoriaSelecionada = filtroCategoria ? filtroCategoria.value : 'todas';
    const servicosFiltrados = categoriaSelecionada === 'todas'
      ? cacheServicosAtuais
      : cacheServicosAtuais.filter(s => s.categoria === categoriaSelecionada);

    renderizarCards(servicosFiltrados, container);
    if (contadorEl) {
      contadorEl.innerText = `${cacheServicosAtuais.length} Serviços Cadastrados`;
    }
  }

  if (filtroCategoria) {
    filtroCategoria.addEventListener('change', () => {
      const categoria = filtroCategoria.value;
      const filtrados = categoria === 'todas'
        ? cacheServicosAtuais
        : cacheServicosAtuais.filter(s => s.categoria === categoria);
      renderizarCards(filtrados, container);
    });
  }

  // Inicialização assíncrona
  await atualizarInterface();

  // Reset do formulário ao abrir modal
  if (modalEl) {
    modalEl.addEventListener('show.bs.modal', () => {
      form.reset();
      form.classList.remove('was-validated');
    });
  }

  // Cadastro de Novo Serviço via HTTP POST
  form.addEventListener('submit', async (e) => {
    e.preventDefault();
    if (!form.checkValidity()) {
      e.stopPropagation();
      form.classList.add('was-validated');
      return;
    }

    const getVal = (id) => document.getElementById(id).value.trim();
    const dados = {
      nome: getVal('nome'),
      categoria: document.getElementById('categoria').value,
      bairro: getVal('bairro'),
      precoBase: parseFloat(getVal('precoBase')) || 0,
      telefone: getVal('telefone'),
      descricao: getVal('descricao')
    };

    if (btnSalvar) btnSalvar.disabled = true;
    if (btnSalvarTexto) btnSalvarTexto.innerText = 'Gravando...';

    try {
      await salvarServico(dados);
      const msg = estaConectadoNaApi()
        ? 'Serviço gravado no banco de dados SQLite!'
        : 'Serviço salvo localmente (Modo Offline).';
      exibirToast(msg, 'success');

      const modal = bootstrap.Modal.getOrCreateInstance(modalEl);
      modal.hide();
      await atualizarInterface();
    } catch (erro) {
      exibirToast('Erro ao processar cadastro.', 'danger');
    } finally {
      if (btnSalvar) btnSalvar.disabled = false;
      if (btnSalvarTexto) btnSalvarTexto.innerText = 'Salvar no Banco';
    }
  });

  // Exclusão Assíncrona via HTTP DELETE
  container.addEventListener('click', async (e) => {
    const btnExcluir = e.target.closest('.btn-excluir');
    if (btnExcluir) {
      const id = btnExcluir.getAttribute('data-id');
      if (confirm('Deseja realmente remover este serviço da vitrine?')) {
        await removerServico(id);
        await atualizarInterface();
        exibirToast('Serviço excluído da vitrine.', 'warning');
      }
    }
  });
});
\end{boxCodigoJS}

% --- FASE 6 ---
\newpage
\begin{boxAtividade}[{Fase 6/6: Verificação de Integração \& Prova de Contingência}]
\textbf{Responsabilidade Arquitetural:} Executar o roteiro de testes em dois cenários: com a API online no laboratório e com o servidor desligado comprovando a resiliência offline.
\end{boxAtividade}

\textbf{1. Teste Conectado (API Online):}
\begin{itemize}[leftmargin=*]
  \item Inicie o backend com \texttt{python app.py} na porta \texttt{5000};
  \item Abra o frontend via Live Server (\texttt{http://127.0.0.1:5500/index.html});
  \item O badge no topo deve exibir: \textbf{\textit{Backend Online (Flask + SQLite)}};
  \item Cadastre um novo serviço e observe a requisição HTTP POST sendo logada no terminal do Flask!
\end{itemize}

\textbf{2. Teste de Resiliência (Fallback LocalStorage):}
\begin{itemize}[leftmargin=*]
  \item No terminal do Flask, pressione \texttt{Ctrl + C} para desligar o servidor;
  \item Recarregue a página no navegador com \texttt{F5};
  \item O frontend detecta a indisponibilidade e muda o badge para: \textbf{\textit{Modo Offline (LocalStorage)}};
  \item Os cards continuam aparecendo normalmente a partir da cópia em cache!
\end{itemize}

\begin{boxTeste}{Critério de Conclusão do Roteiro 05}
A aplicação demonstrou sucesso na comunicação assíncrona cliente-servidor via Fetch API, realizou persistência no SQLite e provou resiliência arquitetural ao operar em modo de contingência offline.
\end{boxTeste}

% =====================================================
% RESUMO DAS TRANSFORMAÇÕES (O QUE MUDOU E POR QUE MUDOU)
% =====================================================

\subsection{Quadro Comparativo: Evolução do Roteiro 04 para o Roteiro 05}

\begin{table}[H]
\centering
\small
\arrayrulecolor{borderGray}
\begin{tabularx}{\textwidth}{|l|X|X|}
\hline
\rowcolor{cinzaCabecalho}
\textbf{Componente} & \textbf{Roteiro 04 (Módulo 2)} & \textbf{Roteiro 05 (Módulo 3)} \\ \hline
\rowcolor{cinzaLinha}
\textbf{Ambiente de Dados} & LocalStorage (isolado na aba do navegador). & Banco relacional SQLite (\texttt{vitrine.db}) no Flask. \\ \hline
\textbf{Comunicação} & Chamadas de função síncronas em memória. & Chamadas assíncronas de rede com \texttt{fetch()} e \texttt{await}. \\ \hline
\rowcolor{cinzaLinha}
\textbf{Segurança de Rede} & Inexistente (tudo no cliente). & Configuração de cabeçalhos CORS (\texttt{flask-cors}). \\ \hline
\textbf{Experiência (UX)} & Resposta imediata sem feedback de rede. & Spinners de espera e badges de status de conectividade. \\ \hline
\rowcolor{cinzaLinha}
\textbf{formatters.js} & Funções puras de formatação. & \textbf{100\% idêntico!} Prova de desacoplamento limpo. \\ \hline
\end{tabularx}
\end{table}

\vspace{0.4cm}
\noindent\textbf{Parabéns!} Você implementou uma arquitetura moderna full-stack desacoplada, apta para o ecossistema profissional e pronta para avançar rumo às Single Page Applications (SPAs) e ao projeto de extensão da SNCT!
